\documentclass[12pt]{article}
\usepackage{amsmath, amssymb}
\usepackage[margin=1in]{geometry}
\setlength{\parskip}{6pt}
\title{QUBO Formulation: Cardinality-Constrained Feature Selection}
\date{}
\begin{document}
\maketitle

\subsection*{Cardinality-Constrained Feature Selection}

\paragraph{Problem Statement.}
Given a set of $N$ candidate features, each associated with a relevance score $s_i$ with respect to a prediction target, and a pairwise redundancy penalty $r_{ij}$ for every distinct pair of features, the task is to select exactly $K$ features. The objective is to maximize the net utility, defined as the sum of the relevance scores of the selected features minus the sum of the redundancy penalties among all selected feature pairs.

\paragraph{Decision Variables.}
For each feature $i \in \{0, 1, \dots, N-1\}$, we introduce a binary decision variable $x_i \in \{0,1\}$. The variable $x_i$ equals $1$ if feature $i$ is selected, and $0$ otherwise. The complete solution is represented by the binary vector $x = (x_0, x_1, \dots, x_{N-1}) \in \{0,1\}^N$.

\paragraph{QUBO Derivation.}
We convert the maximization problem into a minimization problem suitable for QUBO solvers by negating the original objective function.

Step 1 -- Original Objective (Minimization Form):
\begin{align}
    \min_{x \in \{0,1\}^N} \left( -\sum_{i=0}^{N-1} s_i x_i + \sum_{0 \le i < j \le N-1} r_{ij} x_i x_j \right).
\end{align}

Step 2 -- Cardinality Constraint:
The requirement to select exactly $K$ features is expressed as:
\begin{align}
    \sum_{i=0}^{N-1} x_i = K.
\end{align}

Step 3 -- Penalty Term Expansion:
We enforce the constraint via a quadratic penalty with weight $\lambda > 0$:
\begin{align}
    P(x) &= \lambda \left( \sum_{i=0}^{N-1} x_i - K \right)^2 \\
    &= \lambda \left( \sum_{i=0}^{N-1} x_i^2 - 2K \sum_{i=0}^{N-1} x_i + K^2 + \sum_{i \neq j} x_i x_j \right).
\end{align}
Using the idempotent property of binary variables, $x_i^2 = x_i$, and noting that $\sum_{i \neq j} x_i x_j = 2 \sum_{0 \le i < j \le N-1} x_i x_j$, we expand and collect terms:
\begin{align}
    P(x) &= \lambda \left( (1 - 2K) \sum_{i=0}^{N-1} x_i + 2 \sum_{0 \le i < j \le N-1} x_i x_j + K^2 \right) \\
    &= \lambda(1 - 2K) \sum_{i=0}^{N-1} x_i + 2\lambda \sum_{0 \le i < j \le N-1} x_i x_j + \lambda K^2.
\end{align}

Step 4 -- Combined QUBO Hamiltonian:
Adding the negated objective and the penalty term yields the total energy function $H(x)$:
\begin{align}
    H(x) &= \sum_{i=0}^{N-1} \left( -s_i + \lambda(1 - 2K) \right) x_i + \sum_{0 \le i < j \le N-1} \left( r_{ij} + 2\lambda \right) x_i x_j + \lambda K^2.
\end{align}

Step 5 -- Q Matrix Entries and Offset:
The QUBO is defined by a symmetric matrix $Q$ and a constant offset $c$ such that $H(x) = x^\top Q x + c$. Reading off the coefficients from the expanded Hamiltonian gives:
\begin{align}
    Q_{i,i} &= -s_i + \lambda(1 - 2K), \quad \forall i \in \{0, \dots, N-1\}, \\
    Q_{i,j} &= r_{ij} + 2\lambda, \quad \forall 0 \le i < j \le N-1, \\
    Q_{j,i} &= Q_{i,j}, \\
    c &= \lambda K^2.
\end{align}

\paragraph{Penalty Weight Justification.}
The penalty weight $\lambda$ must be chosen sufficiently large to ensure that any violation of the cardinality constraint increases the total energy more than the maximum possible improvement in the original objective. The worst-case gain from flipping a single variable in the original objective is bounded by the sum of all absolute relevance scores and all redundancy penalties. Therefore, setting
\begin{align}
    \lambda > \sum_{i=0}^{N-1} |s_i| + \sum_{0 \le i < j \le N-1} r_{ij}
\end{align}
guarantees that the penalty term strictly dominates any feasible objective improvement, ensuring that infeasible configurations are energetically unfavorable.

\paragraph{Correctness Argument.}
Minimizing $H(x)$ over $\{0,1\}^N$ recovers the optimal feature selection because (i) any assignment satisfying $\sum_{i=0}^{N-1} x_i = K$ incurs zero penalty, so its energy exactly equals the negated original objective value; (ii) any infeasible assignment violates the cardinality constraint by at least one unit, incurring a penalty strictly greater than the maximum possible reduction in the original objective. Consequently, the energy of any infeasible solution exceeds that of at least one feasible solution. The global minimum of $H(x)$ must therefore occur at a feasible assignment that minimizes the negated objective, which is equivalent to maximizing the original net utility.

\paragraph{Illustrative Example.}
Consider an instance with $N=3$ features and $K=2$. Let the relevance scores be $s_0 = 10$, $s_1 = 8$, $s_2 = 6$, and the redundancy penalties be $r_{01} = 1$, $r_{02} = 2$, $r_{12} = 3$. We choose $\lambda = 100$, which satisfies the lower bound $\sum |s_i| + \sum r_{ij} = 24 + 6 = 30$.

Instantiating the general formulas yields the following Q matrix entries and offset:
\begin{align}
    Q_{0,0} &= -10 + 100(1 - 4) = -310, \\
    Q_{1,1} &= -8 + 100(1 - 4) = -298, \\
    Q_{2,2} &= -6 + 100(1 - 4) = -276, \\
    Q_{0,1} &= 1 + 200 = 201, \\
    Q_{0,2} &= 2 + 200 = 202, \\
    Q_{1,2} &= 3 + 200 = 203, \\
    c &= 100 \times 2^2 = 400.
\end{align}
The concrete Hamiltonian is:
\begin{align}
    H(x) &= -310 x_0 - 298 x_1 - 276 x_2 + 201 x_0 x_1 + 202 x_0 x_2 + 203 x_1 x_2 + 400.
\end{align}
Since $K=2$, there are $\binom{3}{2}=3$ feasible assignments. Evaluating $H(x)$ for each:
\begin{itemize}
    \item $x = (1,1,0)$: $H = -310 - 298 + 201 + 400 = -17$.
    \item $x = (1,0,1)$: $H = -310 - 276 + 202 + 400 = -14$.
    \item $x = (0,1,1)$: $H = -298 - 276 + 203 + 400 = -11$.
\end{itemize}
The global minimum is $H(x^*) = -17$ at $x^* = (1,1,0)$, corresponding to selecting features $0$ and $1$. The maximum total score is $10 + 8 - 1 = 17$, matching the negated energy $-H(x^*) + c = 17$.
\end{document}